\section{Experiments and Results}
\label{sec:experiments}

We report eight controlled experiments grouped into two regimes.
\Cref{sec:experiments-de} (Runs~1--5) uses the Dynamic Equilibrium
specialisation of EDEM with linear supply and demand schedules,
discrete sales, and a rolling-mean clearing price; these isolate the
effects of estimation noise, patience, agent density, and shocks
under a textbook population balancer. \Cref{sec:experiments-edem}
(Runs~6--8) uses the speculative-market EDEM variant with
multiplicative value updates; these illustrate the bubble, balancer
sign, and rising-error regimes that the static
divergence-of-opinion literature cannot reach.

Unless stated otherwise, all runs use a $32 \times 32$ toroidal grid,
balance period $T_{B} = 100$ ticks, sale window $W = 25$ sales, and
the bid-acceptance rule \cref{eq:cond2}. Each experiment is replicated
across at least eight independent seeds; figures plot the median across
seeds plus the 10th--90th percentile band, with individual seed
trajectories overlaid faintly. Full parameter tables are in
\cref{app:params}.

% =====================================================================
\subsection{Dynamic Equilibrium experiments}
\label{sec:experiments-de}
% =====================================================================

\subsubsection{Run 1 -- Stable equilibrium under favourable conditions}
\label{sec:exp-stable}

We begin from the textbook equilibrium $p^{*} = 100$, $q^{*} = 50$
defined by $\hat{Q}_{s}(p) = 0.5\,p$ and
$\hat{Q}_{d}(p) = 100 - 0.5\,p$, with maximum estimation error
$\bar{\sigma} = 5$ percent, patience drawn uniformly on $[0, 50]$,
and balance period $T_{B} = 100$. \Cref{fig:run1-stable} shows the
median market price hovering close to $p^{*}$: across ten seeds and
20{,}000 ticks, the maximum upward deviation is $+10.7\%$ and the
maximum downward deviation $-2.4\%$. The companion buyer/seller
counts oscillate within $\pm 5$ of the equilibrium quantity. The
slight upward bias of the median trajectory is consistent with
\citet{arbuzov2018de}'s reported $-4.9\%/+6.2\%$ envelope at
100{,}000 ticks; we attribute it to the bid-selection asymmetry
discussed in \cref{sec:framework-de-special}, which the rolling
average partially absorbs.

\begin{figure}[t]
\centering
\includegraphics[width=\linewidth]{fig1_stable}
\caption{Run~1 -- Stable equilibrium with low valuation error and
moderate patience. Top: market price (median, with 10--90 percentile
band; faint individual-seed traces). Bottom: agent counts. Dashed
line: textbook equilibrium price.}
\label{fig:run1-stable}
\end{figure}

\subsubsection{Run 2 -- Imprecise valuations cause business cycles}
\label{sec:exp-cycles}

Holding all other parameters at Run~1 values and increasing
$\bar{\sigma}$ from $5$ to $25$ percent, the market loses its
narrow-band stability. \Cref{fig:run2-cycles} shows recurrent
oscillations of the market price between roughly $+34\%$ and
$-17\%$ of equilibrium, accompanied by counter-phased oscillations
of buyer and seller populations: episodes of overbidding draw
sellers into the market while buyers retreat, after which the
seller surplus pushes prices back down and the cycle reverses.
This is a pure consequence of estimation noise --- there are no
exogenous shocks --- and it suggests that observed business cycles
in real markets need not require any external driver beyond
heterogeneous misvaluation, provided the divergence of opinion is
large relative to the balancer's adjustment speed.

\begin{figure}[t]
\centering
\includegraphics[width=\linewidth]{fig2_high_error}
\caption{Run~2 -- Imprecise valuations (\(\bar{\sigma} = 25\%\))
generate endogenous business cycles. Same conventions as
\cref{fig:run1-stable}.}
\label{fig:run2-cycles}
\end{figure}

\subsubsection{Run 3 -- Patience pushes price above equilibrium}
\label{sec:exp-patience}

Doubling the maximum seller patience from $50$ to $100$ ticks
(holding $\bar{\sigma}$ at the Run~1 value of $5\%$) produces a
durable upward shift in the market price.
\Cref{fig:run3-patience} shows the median price climbing
monotonically to roughly $+14\%$ above $p^{*}$ over 20{,}000 ticks
and showing no sign of reverting. The seller population settles
at $\sim 57$ (about $+14\%$ of $q^{*}$) and the buyer population
at $\sim 41$ ($-18\%$). Patience increases sellers' time-on-market
and thus their willingness to wait for richer bids; the realised
sale prices tilt up; the rolling average follows.
\citet{arbuzov2018de} reported a similar
$+10\%/+10\%/-20\%$ shift, attributing it to seller forbearance
on stagnating real-estate markets where neither side can quickly
exit.

\begin{figure}[t]
\centering
\includegraphics[width=\linewidth]{fig3_patience}
\caption{Run~3 -- Doubling seller patience to $100$ ticks pushes the
market price about $14\%$ above the textbook equilibrium and
sustains it there. Same conventions as \cref{fig:run1-stable}.}
\label{fig:run3-patience}
\end{figure}

\subsubsection{Run 4 -- Low agent density pushes price below equilibrium}
\label{sec:exp-density}

Shifting the demand intercept downward by $50$ units yields a new
textbook equilibrium $(p^{*}, q^{*}) = (50, 25)$. With only $\sim
25$ buyers and $\sim 25$ sellers on a $32 \times 32$ grid, encounters
are sparse: a typical buyer's wiggle-and-step trajectory rarely
crosses a seller's home before the seller's patience timer expires.
Sellers therefore lower ask prices repeatedly without receiving
adequate bids, and the realised sale prices --- and hence the rolling
market price --- settle \emph{below} the textbook equilibrium.
\Cref{fig:run4-density} shows a steady drift to roughly $p \approx
37$, an $18\%$ shortfall, with the seller population shrinking to
$\sim 19$ and the buyer population growing to $\sim 30$ as the
balancer chases the depressed price. Patience and density are partial
substitutes (\cref{sec:bg-dynamics}): both proxy for time-on-market.

\begin{figure}[t]
\centering
\includegraphics[width=\linewidth]{fig4_low_density}
\caption{Run~4 -- Halving the demand intercept reduces equilibrium
quantity to $25$. Sparse encounters cause sellers to lower ask prices
repeatedly; the realised market price settles about $18\%$ below the
textbook equilibrium. Same conventions as \cref{fig:run1-stable}.}
\label{fig:run4-density}
\end{figure}

\subsubsection{Run 5 -- Shock handling and transitional markets}
\label{sec:exp-shock}

The previous four runs evolve under fixed parameters. Run~5
introduces shocks via the
\texttt{set\_demand}/\texttt{set\_supply} hooks
(\cref{sec:implementation-pseudocode}) and demonstrates how the model
absorbs them. \Cref{fig:run5-shock} reports two scenarios.

\paragraph{Scenario A (single shock).} The market begins in the Run~1
equilibrium. At $t = 3000$ the demand intercept drops to $50$,
collapsing the textbook equilibrium to $p^{*} = 50$. The realised
price descends slowly --- it has not finished adjusting by $t = 7000$
when seller patience is raised to $165$. The combination of a
lowered demand and elevated patience continues to drive the market
downward (because the price is still adjusting from the previous
$p^{*} = 100$), and even by $t = 12000$ the realised price has not
settled. Doc 2's claim that high patience would lift the price
\emph{back to} $p^{*} = 50$ presumes a fully-adjusted starting state
that is not reached within twelve thousand ticks; the figure makes
visible the more general \citet{arbuzov2018de} observation that
``adjustment to single shock takes a long time.''

\paragraph{Scenario B (transitional market).} The demand intercept
toggles between $125$ and $75$ every $2000$ ticks. The textbook
equilibrium accordingly steps between $p^{*} = 125$ and $p^{*} = 75$,
but the realised market price never settles at either: each shock
is absorbed only partially before the next arrives. The figure
confirms a market \emph{in constant transition from one unknown
state to another} --- the empirical foothold for
\cref{sec:framework-schedules}'s claim that supply and demand are
sample paths of a stochastic process, not deterministic schedules.

\begin{figure}[t]
\centering
\includegraphics[width=\linewidth]{fig5_shock}
\caption{Run~5 -- Two shock scenarios. Top: a single demand shock at
$t = 3000$ followed by a patience boost at $t = 7000$; the market
absorbs the shock slowly and has not fully adjusted by the end of
the simulation. Bottom: a transitional market in which the demand
intercept toggles every $2000$ ticks. The market price never
converges to either textbook equilibrium. Red dashed verticals mark
shock events.}
\label{fig:run5-shock}
\end{figure}

% =====================================================================
\subsection{Speculative-market EDEM experiments}
\label{sec:experiments-edem}
% =====================================================================

The remaining three runs use the EDEM variant of
\cref{sec:framework-de-special} with multiplicative value updates
\cref{eq:value-update}. We track the dimensionless ratio
$v_{t}(h) / v^{*}(h)$ averaged over all homes; a value of $1$
indicates fair valuation, values above $1$ indicate the market has
priced homes above their unobservable fair value.

\subsubsection{Run 6 -- Speculative bubble at $C_{b} = 0$}
\label{sec:exp-bubble}

With balancer disabled ($C_{b} = 0$) and divergence of opinion
$\bar{\sigma} = 15\%$, the market exhibits a clean log-linear bubble.
\Cref{fig:run6-bubble} shows the median value-to-true ratio
multiplying by $\sim 421\times$ over $3000$ ticks ($150$ epochs),
with the band tightly enclosing the median: the bubble is
deterministic in shape and stochastic only in slope. Buyer and
seller populations are frozen at $20$ each, since $C_{b} = 0$
admits no balancer adjustment. The mechanism is the bid-selection
asymmetry of \cref{sec:framework-de-special}: each epoch's average
winning ratio $\bar{r}_{t}$ is $> 1$ in expectation because winning
bids are the maxima of multiple noisy draws, and \cref{eq:value-update}
compounds this multiplicatively over epochs.

\begin{figure}[t]
\centering
\includegraphics[width=\linewidth]{fig6_bubble}
\caption{Run~6 -- Speculative bubble at $C_{b} = 0$,
$\bar{\sigma} = 15\%$. Top: dimensionless market value $v_{t}/v^{*}$
on a logarithmic axis; the trajectory is approximately linear in
$\log v$, indicating exponential growth. Bottom: agent populations
remain at $(20, 20)$.}
\label{fig:run6-bubble}
\end{figure}

\subsubsection{Run 7 -- Balancer sign shapes (but cannot eliminate)
                       the bubble}
\label{sec:exp-balancer-sweep}

Holding $\bar{\sigma}$ at $15\%$ and varying $C_{b} \in \{+1, 0, -1\}$
yields three qualitatively distinct trajectories overlaid in
\cref{fig:run7-balancer-sweep}. The mean-reverting balancer
($C_{b} = +1$) dampens the bubble most aggressively: the value ratio
plateaus near $4 \times$ as the seller population grows to roughly
$40$ and the buyer population shrinks to one. The trend-following
balancer ($C_{b} = -1$) admits a larger bubble than $+1$, ending
near $7.7\times$, with mirrored populations (one seller, $\sim 40$
buyers). The no-balancer case ($C_{b} = 0$) dominates both, ending
at $\sim 421\times$.

This last observation refines the framing in
\citet{arbuzov2018edem}, which described the trend-following case
as the most explosive (the ``Silicon Valley real estate'' or
``bitcoin'' regime). In the Mesa simulation the most explosive case
is the no-balancer one: removing one population pole entirely (as
trend-following does) reduces the multiplicative pressure on each
remaining seller, because there are simply fewer sellers whose
values can be inflated. The trend-following case is more aggressive
than mean-reverting, but neither is as aggressive as no balancer.

The companion finding is that no setting of $C_{b}$ in this
parameter grid restores $v_{t} / v^{*} = 1$. Bid-selection
asymmetry introduces a positive bias that even a strong
mean-reverting balancer cannot fully cancel within the timescales
considered.

\begin{figure}[t]
\centering
\includegraphics[width=\linewidth]{fig7_balancer_sweep}
\caption{Run~7 -- Balancer-coefficient sweep at $\bar{\sigma} = 15\%$.
Three regimes are overlaid: mean-reverting ($C_{b} = +1$), no
balancer ($C_{b} = 0$), and trend-following ($C_{b} = -1$). The
no-balancer case dominates; the trend-following case beats the
mean-reverting case but is itself bounded above by the no-balancer
case. Bottom panel shows the population swings induced by each
balancer.}
\label{fig:run7-balancer-sweep}
\end{figure}

\subsubsection{Run 8 -- Trend-following balancer with rising error}
\label{sec:exp-double-exp}

Doc 2's stylised ``bitcoin'' or ``Silicon Valley'' regime combines
trend-following ($C_{b} = -1$) with a divergence of opinion that
itself grows over time, modelling the historical observation that
late-stage bull markets see widening dispersion in analyst forecasts.
We initialise $\bar{\sigma}_{0} = 0.05$ (i.e.\ $5\%$) and let
$\bar{\sigma}_{t+T} = \bar{\sigma}_{t} + 0.005$ each epoch
(equivalently, $+0.5$ percentage points), so $\bar{\sigma}$ reaches
$0.80$ ($80\%$) by the end of the run.
\Cref{fig:run8-double-exp} shows the resulting trajectory: the
value ratio climbs to $\sim 30 \times$ over $3000$ ticks. The
trajectory is super-linear in linear time but visibly concave-down
in log-time, indicating sustained exponential rather than
double-exponential growth at this parameter setting; we conjecture
that genuine double-exponential acceleration would emerge for
larger $\bar{\sigma}$ growth rates than we explore here. The
population panel mirrors Run~7's $C_{b} = -1$ behaviour: sellers
drained to one, buyers saturated at $\sim 44$.

\begin{figure}[t]
\centering
\includegraphics[width=\linewidth]{fig8_double_exp}
\caption{Run~8 -- Trend-following balancer ($C_{b} = -1$) combined
with linearly-growing divergence of opinion
($\bar{\sigma}: 5 \to 80$). Top panel: market value on a log scale.
Middle: agent populations. Bottom: $\bar{\sigma}$ over time. The
trajectory is sustained exponential rather than visibly
double-exponential.}
\label{fig:run8-double-exp}
\end{figure}

\subsubsection{Run 9 -- Sensitivity grid}
\label{sec:exp-sensitivity}

The eight runs reported above each fix a single parameter vector. To
rule out the possibility that the order-statistic drift of
\cref{sec:framework-asymmetry} is an artefact of one well-chosen
calibration, we sweep $C_{b}$ and $\bar{\sigma}$ over a $5 \times 6$
grid (\(C_{b} \in \{-1, -0.5, 0, +0.5, +1\}\),
$\bar{\sigma} \in \{5, 10, 15, 20, 25, 30\}\%$), running five seeds per
cell for $1500$ ticks. \Cref{fig:run9-sensitivity} reports the
median terminal $v_{t}/v^{*}$ across seeds, on a $\log_{10}$ colour
scale.

Three observations stand out:

\begin{enumerate}[leftmargin=*,nosep]
\item \textbf{The drift is universal in this range.} Every cell of the
grid has $v_{t}/v^{*} > 1$, and 100\% of cells exceed
$1.5 \times$. The order-statistic mechanism is not a corner-case.
\item \textbf{Bubble territory is large.} 43\% of cells exceed
$10 \times$ within $1500$ ticks; growth at $\bar{\sigma} = 30\%$
reaches $235 \times$ in the no-balancer column.
\item \textbf{$C_{b} = 0$ dominates within each $\bar{\sigma}$ column.}
This confirms the Run~7 finding (\cref{sec:exp-balancer-sweep}) at
five distinct $\bar{\sigma}$ values: removing the balancer is more
explosive than either mean-reverting or trend-following balancers,
across the whole grid.
\end{enumerate}

The grid does not constitute a closed-form proof that the drift
survives \emph{all} parameter regimes; in particular, very large
$|C_{b}|$ values that drain one side of the population to a single
agent within a few epochs may behave differently. Section
\ref{sec:disc-balancer} discusses the limits of the linear-balancer
intervention more carefully.

\begin{figure}[t]
\centering
\includegraphics[width=\linewidth]{fig9_sensitivity}
\caption{Run~9 -- Sensitivity grid: median terminal $v_{t}/v^{*}$
over a $5 \times 6$ sweep of balancer coefficient $C_{b}$ versus
divergence of opinion $\bar{\sigma}$, with 5 seeds per cell at
$t = 1500$. Colour: $\log_{10}$ of the cell median; cell text:
linear value. Red squares mark the cells corresponding to Runs~6
and~7 ($C_{b} \in \{-1, 0, +1\}$ at $\bar{\sigma} = 15\%$). The
order-statistic drift produces $v_{t}/v^{*} > 1$ in every cell.}
\label{fig:run9-sensitivity}
\end{figure}

% =====================================================================
\subsection{Summary across runs}
\label{sec:exp-summary}
% =====================================================================

\Cref{tab:experiment-summary} consolidates the end-of-run statistics.
The eight runs sweep three parameter axes: estimation noise
($\bar{\sigma}$), seller patience, and agent density (Runs 1--4); a
shock schedule (Run~5); and the EDEM balancer parameter $C_{b}$ with
a fixed or growing $\bar{\sigma}$ (Runs 6--8). Together they delineate
six qualitatively distinct steady-state regimes -- band-stable,
business-cycle, persistent overshoot, persistent undershoot, runaway
exponential, and constant transition -- all reachable from the same
underlying agent ruleset.

\begin{table}[t]
\centering
\small
\caption{Eight EDEM runs: parameters and observed outcomes
at the end of each simulation. Median across $\ge 8$ seeds.
DE runs report median deviation of the rolling market price from the
textbook equilibrium; EDEM runs report the dimensionless
$v_{t}/v^{*}$ ratio.}
\label{tab:experiment-summary}
\begin{tabular}{l l l l}
\toprule
\textbf{Run} & \textbf{Manipulated parameter} & \textbf{Outcome metric} & \textbf{Regime} \\
\midrule
1 & $\bar{\sigma} = 5\%$, patience $\le 50$       & $-2\%$ to $+11\%$ band       & band-stable        \\
2 & $\bar{\sigma} = 25\%$                         & $-17\%$ to $+34\%$, periodic & business-cycle     \\
3 & patience $\le 100$                            & $+14\%$ persistent           & overshoot          \\
4 & demand intercept $-50$                        & $-18\%$ persistent           & undershoot         \\
5 & shock at $t=3000$, then re-shocks             & price chases moving target   & transitional       \\
\midrule
6 & EDEM, $C_{b} = 0$, $\bar{\sigma}=15\%$        & $v_{t}/v^{*} = 421\times$    & runaway bubble     \\
7 & EDEM, $C_{b} \in \{-1, 0, +1\}$               & $\{7.7, 421, 4.2\}\times$    & balancer-shaped    \\
8 & EDEM, $C_{b} = -1$, $\bar{\sigma}: 5\to80$    & $v_{t}/v^{*} = 30\times$     & sustained super-linear \\
\bottomrule
\end{tabular}
\end{table}
